\begin{figure}[!htbp]
\centering
\begin{tikzpicture}[
  font=\small,
  boundary/.style={draw=black!85, line width=.8pt},
  orient/.style={-{Stealth[length=4pt]}, draw=black!85, line width=.8pt},
  mesh/.style={draw=black!25, line width=.35pt}
]
\node at (8.02,4.25)
  {$\displaystyle\Phi(\sigma,\vartheta)
    =\operatorname{Fl}^{1}_{\sigma f_1+\vartheta f_2}(p_T),
    \qquad f_i=a_pE_i$};
\node[font=\small\bfseries] at (2.28,3.55) {Parameter square};
\node[font=\small\bfseries] at (12.52,3.55) {Image surface (schematic)};

\coordinate (o) at (.80,.30);
\coordinate (a) at (3.50,.30);
\coordinate (b) at (3.50,3.00);
\coordinate (c) at (.80,3.00);
\fill[black!3] (o) rectangle (b);
\draw[mesh] (1.70,.30) -- (1.70,3.00);
\draw[mesh] (2.60,.30) -- (2.60,3.00);
\draw[mesh] (.80,1.20) -- (3.50,1.20);
\draw[mesh] (.80,2.10) -- (3.50,2.10);
\draw[boundary] (o) -- (a) -- (b) -- (c) -- cycle;
\draw[orient] (o) -- (2.32,.30);
\draw[orient] (a) -- (3.50,1.82);
\draw[orient] (b) -- (1.98,3.00);
\draw[orient] (c) -- (.80,1.48);
\fill (o) circle (1.6pt);
\node[anchor=north east] at (o) {$(0,0)$};
\node[anchor=north west] at (a) {$(h,0)$};
\node[anchor=south west] at (b) {$(h,h)$};
\node[anchor=south east] at (c) {$(0,h)$};
\node[fill=white, inner sep=3pt] at (2.15,1.65) {$[0,h]^2$};

\draw[-{Stealth[length=6pt]}, thick] (4.63,1.77) -- (8.65,1.77)
  node[midway, above=5pt] {$\Phi$}
  node[midway, below=6pt, align=center, text width=3.85cm]
    {time-one flow of\\the field\\
     $\sigma f_1+\vartheta f_2$};

\coordinate (p) at (9.72,.37);
\coordinate (q) at (14.40,.65);
\coordinate (r) at (14.67,2.85);
\coordinate (s) at (10.17,3.00);
\path[fill=black!3]
  (p) .. controls (11.15,.80) and (12.85,.10) .. (q)
      .. controls (14.13,1.31) and (15.18,2.03) .. (r)
      .. controls (13.28,3.27) and (11.75,2.53) .. (s)
      .. controls (10.62,2.15) and (9.37,1.32) .. cycle;
\draw[mesh] (11.23,.56)
  .. controls (11.02,1.25) and (11.88,2.24) .. (11.64,2.84);
\draw[mesh] (12.84,.40)
  .. controls (12.58,1.30) and (13.51,2.24) .. (13.17,3.01);
\draw[mesh] (9.72,1.20)
  .. controls (11.20,1.72) and (13.01,.84) .. (14.45,1.38);
\draw[mesh] (10.09,2.12)
  .. controls (11.68,2.52) and (13.09,1.71) .. (14.83,2.11);
\draw[orient] (p) .. controls (11.15,.80) and (12.85,.10) .. (q);
\draw[orient] (q) .. controls (14.13,1.31) and (15.18,2.03) .. (r);
\draw[orient] (r) .. controls (13.28,3.27) and (11.75,2.53) .. (s);
\draw[orient] (s) .. controls (10.62,2.15) and (9.37,1.32) .. (p);
\fill (p) circle (1.8pt);
\node[anchor=north, inner sep=5pt] at (10.26,.37) {$p_T=\Phi(0,0)$};
\node[fill=white, inner sep=3pt] at (12.28,1.72) {$\Phi([0,h]^2)$};
\node[anchor=north] at (12.51,-.24) {$\ell_h=\Phi(\partial[0,h]^2)$};

\node[align=center] at (8.02,-1.26)
  {$\displaystyle\frac12\mathcal H_T''(0)
     =\pm a_p^2R^{g'}(E_1,E_2)(p_T),
   \qquad \mathcal L_T(h)\le a_phQ_T,
   \quad Q_T=e^{o(T)}$.};
\end{tikzpicture}
\caption{The loop surface of Equation~\eqref{holo:surface}, written on
$[0,h]^2$ rather than the unit square. Each image point is the time-one
flow of the single vector field $\sigma f_1+\vartheta f_2$ from $p_T$;
the oriented boundary gives the based loop. The drawing is schematic,
with $a_p$ constant on the observation box. Lemma~\ref{holo:loops}
gives the quadratic holonomy coefficient (up to the stated sign
convention) and the developed length in the transported test metric.}
\label{fig:parameter-flow-loop}
\end{figure}
